\subsection{Оператор с ядром Гильберта—Шмидта}

Будем рассматривать уравнение второго рода

\stepcounter{equation}
\begin{equation}\tag{1}\label{eq:hs-eq}
  u(x) - \int\limits_{0}^{\ell} K(x,y)\, u(y)\, dy = f(x), \qquad x \in D ,
\end{equation}
и всюду дальше полагать

\stepcounter{equation}
\begin{equation}\tag{2}\label{eq:hs-kernel}
  \int\limits_{0}^{\ell}\int\limits_{0}^{\ell} \abs{K(x,y)}^2\, dy\, dx < \infty
  \qquad\Longleftrightarrow\qquad
  K \in \Lp{2}(D \times D) .
\end{equation}
Такое ядро называется ядром \textbf{Гильберта—Шмидта}.

\begin{definition}
Интегральный оператор с ядром $K$:
\[
  (Au)(x) = \int\limits_{0}^{\ell} K(x,y)\, u(y)\, dy .
\]
\end{definition}

\begin{theorem}\label{thm:hs-compact}
Ядро Гильберта—Шмидта определяет компактный (вполне непрерывный)
интегральный оператор в пространстве $\Lp{2}(D)$, норма которого
\[
  \norm{A} \le \left(\ \int\limits_{0}^{\ell}\int\limits_{0}^{\ell}
  \abs{K(x,y)}^2\, dy\, dx \right)^{1/2} .
\]
\end{theorem}

\subsubsection*{Шаг 1: ограниченность}

По теореме Фубини из~\eqref{eq:hs-kernel} следует, что для почти всех $x$
выполнено $K(x,\cdot) \in \Lp{2}(D)$. Тогда для почти всех $x$ по неравенству
Коши—Буняковского
\[
  \abs{(Au)(x)}^2
  \le \left[\ \int\limits_{0}^{\ell} \abs{K(x,y)\, u(y)}\, dy \right]^2
  \le \int\limits_{0}^{\ell} \abs{K(x,y)}^2 dy \cdot
      \underbrace{\int\limits_{0}^{\ell} \abs{u(y)}^2 dy}_{=\ \norm{u}^2} .
\]
Интегрируя по $x$:
\[
  \underbrace{\int\limits_{0}^{\ell} \abs{(Au)(x)}^2 dx}_{=\ \norm{Au}^2}
  \ \le\ \int\limits_{0}^{\ell}\int\limits_{0}^{\ell} \abs{K(x,y)}^2\, dy\, dx
  \cdot \norm{u}^2 ,
\]
то есть

\stepcounter{equation}
\begin{equation}\tag{3}\label{eq:hs-bound}
  \norm{Au} \le \norm{K}\cdot\norm{u} ,
  \qquad
  \norm{K} := \left(\iint\limits_{D\times D}\abs{K(x,y)}^2 dy\,dx\right)^{1/2} .
\end{equation}
Итого $A$ — ограниченный оператор, и $\norm{A} \le \norm{K}$.

\subsubsection*{Шаг 2: разложение ядра по базису}

Пространство $\Lp{2}(D)$ сепарабельно, значит в нём существует
ортонормированный базис $\{\varphi_k\}_{k=1}^{\infty}$. Тогда (по известной
лемме) система произведений
\[
  \big\{\, \varphi_k(x)\,\overline{\varphi_i(y)} \,\big\}_{k,i=1}^{\infty}
\]
образует ортонормированный базис в $\Lp{2}(D \times D)$. Поэтому ядро
раскладывается в сходящийся в $\Lp{2}(D\times D)$ ряд
\[
  K(x,y) = \sum_{k,i=1}^{\infty} k_{ki}\, \varphi_k(x)\,\overline{\varphi_i(y)} ,
  \qquad
  \sum_{k,i=1}^{\infty} \abs{k_{ki}}^2 = \norm{K}^2 < \infty .
\]

Возьмём частичную сумму
\[
  K_N(x,y) = \sum_{k,i=1}^{N} k_{ki}\, \varphi_k(x)\,\overline{\varphi_i(y)} .
\]

\subsubsection*{Шаг 3: обрезанный оператор конечномерен}

Для оператора $A_N$ с ядром $K_N$
\[
  (A_N u)(x)
  = \int\limits_{D} \sum_{k,i=1}^{N} k_{ki}\,\varphi_k(x)\,\overline{\varphi_i(y)}\,
    u(y)\, dy
  = \sum_{k,i=1}^{N} k_{ki}\,\varphi_k(x)
    \int\limits_{D} u(y)\,\overline{\varphi_i(y)}\, dy .
\]
Значения $A_N u$ при любом $u$ лежат в линейной оболочке
$\varphi_1,\dots,\varphi_N$, то есть $A_N$ — \textbf{конечномерный} оператор
(оператор конечного ранга). Такой оператор компактен: образ ограниченного
множества лежит в ограниченном подмножестве конечномерного пространства, а
там работает теорема Больцано—Вейерштрасса.

\subsubsection*{Шаг 4: предельный переход}

Применим оценку~\eqref{eq:hs-bound} к оператору $A - A_N$, ядро которого есть
$K - K_N$:
\[
  \norm{A - A_N}
  \ \le\ \left[\ \iint\limits_{D \times D}
  \abs{K(x,y) - K_N(x,y)}^2\, dy\, dx \right]^{1/2}
  = \Bigg(\sum_{\substack{k,i \\ \max(k,i) > N}} \abs{k_{ki}}^2\Bigg)^{1/2}
  \xrightarrow[N \to \infty]{} 0
\]
как остаток сходящегося ряда. Значит $A$ — предел по операторной норме
последовательности компактных операторов, а потому сам компактен.

\begin{idea}[title={Почему предел компактных операторов компактен}]
Это стандартный факт, но выкладка короткая. Пусть $\norm{u_n} \le M$. Оператор
$A_1$ компактен, поэтому из $\{u_n\}$ выбирается подпоследовательность, на
которой $\{A_1 u_n\}$ сходится; из неё — подпоследовательность для $A_2$, и так
далее. Диагональная последовательность $g_n$ годится сразу для всех $A_N$.
Тогда
\[
  \norm{A g_n - A g_p}
  \le \underbrace{\norm{(A - A_N) g_n}}_{\le\, M\norm{A-A_N}}
    + \norm{A_N g_n - A_N g_p}
    + \underbrace{\norm{(A_N - A) g_p}}_{\le\, M\norm{A-A_N}} .
\]
Сначала выбираем $N$ так, чтобы крайние слагаемые были малы, затем — номер,
начиная с которого мал средний член. Получаем фундаментальность, а $\Lp{2}(D)$
полно.
\end{idea}

\begin{idea}[title={Второй способ: приближение непрерывным ядром}]
Тот же результат получается и без разложения по базису. Непрерывные функции
плотны в $\Lp{2}(D\times D)$, поэтому ядро приближается непрерывным $K_\eps$ с
$\norm{K - K_\eps} < \eps$. Для непрерывного ядра на компакте семейство
образов ограниченного множества равномерно ограничено и равностепенно
непрерывно (обе оценки — то же неравенство Коши—Буняковского), так что
компактность даёт теорема Арцела—Асколи. Дальше — тот же предельный переход
шага~4.

\medskip
Вариант с базисом предпочтительнее потому, что попутно выясняется главное:
\textbf{ядро Гильберта—Шмидта сколь угодно точно приближается вырожденными
ядрами}. Именно поэтому следующий шаг теории — разобрать уравнение с
вырожденным ядром до конца (\S4): общий случай получится предельным
переходом.
\end{idea}
